% !TEX root = ../main.tex
\section{Q01: De 5 centertyper -- fordele, ulemper og forutsetninger}\label{q:01}

\begin{quote}
\emph{I relation til målinger af organisatoriske enheders præstationer bedes du redegøre for de 5 centertyper. I forlængelse heraf ønskes fordele og ulemper diskuteret, ligesom der ønskes argumenteret for, under hvilke forudsætninger de enkelte centertyper bør vælges.}
\end{quote}

\subsection*{Svar (kort)}
De fem centertypene -- \emph{cost, expense, revenue, profit, investment} -- representerer en stige av delegert beslutningsrett, og valget av centertype må følge controllability-prinsippet: måltallet skal dekke akkurat det lederen faktisk kan kontrollere, verken mer eller mindre. Centertypen er derfor ikke et isolert valg av prestasjonsmål (S2), men et samtidig design av beslutningsrettigheter (S1), prestasjonsmåling (S2) og insentiver (S3). Fordelene og ulempene er en direkte konsekvens av hvor godt centertypens måltall stemmer med disse tre dimensjonene. [SLIDE]

\subsection*{Relevante teorier (kort)}

\paragraph{\thref{responsibility-centers}{Responsibility centers -- de 5 centertyper}\quad\SOne/\STwo/\SThree}
Cost center måler kostnad ved gitt output (produksjon, sykehus). Expense center måler budsjettoverholdelse når output er vanskelig å måle (HR, R\&D). Revenue center måler omsetning når salgsinnsats kontrolleres lokalt men pris ikke (salgsavdelinger). Profit center måler profit når lederen styrer både kostnad og pris (divisjoner). Investment center måler avkastning på kapital (ROA, EVA) når også investeringer er delegert. Hver type kobler beslutningsrett (S1) til måltall (S2) og videre til belønning (S3). [SLIDE]

\paragraph{\thref{controllability-principle}{Controllability-prinsippet}\quad\STwo}
Lederen skal kun måles på utfall hun selv kontrollerer; ikke-kontrollerbar støy må filtreres bort eller kompenseres med høyere fastlønn. Dette er det normative argumentet for hvorfor centertypene må skreddersys: en produksjonssjef uten prisansvar skal ikke evalueres på profit, og en salgssjef uten kostnadsansvar skal ikke straffes for høye produksjonskostnader. Prinsippet leder direkte til valgkriteriene for hvilken centertype som passer. [SLIDE/BOK]

\subsection*{Illustrasjon}

\begin{center}
\renewcommand{\arraystretch}{1.15}
\begin{tabular}{@{}lccccl@{}}
\toprule
\textbf{Centertype} & \textbf{Kost.} & \textbf{Innt.} & \textbf{Profit} & \textbf{Kap.} & \textbf{Måltall} \\
\midrule
Cost center & \checkmark & & & & Kostnad pr.\ enhet \\
Expense center & (budsj.) & & & & Budsjettoverholdelse \\
Revenue center & & \checkmark & & & Omsetning \\
Profit center & \checkmark & \checkmark & \checkmark & & Regnskapsmessig profit \\
Investment center & \checkmark & \checkmark & \checkmark & \checkmark & ROA / EVA / IRR \\
\bottomrule
\end{tabular}
\end{center}
\textit{Tabell: De fem centertypene som stige av delegerte beslutningsrettigheter og tilhørende måltall. [SLIDE -- BSZ kap.\ 17]}

% Slide-figur fra BSZ kap. 17, slide 4 finnes som EMF og må konverteres til PDF/PNG for å \includegraphics-es i denne notatsamlingen.
% Når konvertert til BSZ_kap_17/by_slide/slide_004/BSZ_kap_17_slide_004_asset_01_image1.pdf, aktiveres:
% \begin{figure}[h]
%   \centering
%   \includegraphics[width=0.7\linewidth]{BSZ_kap_17/by_slide/slide_004/BSZ_kap_17_slide_004_asset_01_image1.pdf}
%   \caption{De fem centertypene -- forelesningsslide. [SLIDE]}
% \end{figure}

\subsection*{Fordele og ulemper -- kort oppsummert}
\begin{itemize}\setlength{\itemsep}{2pt}
\item \textbf{Cost center.} Fordel: enkelt å måle, klart ansvar. Ulempe: kostnadsminimering $\neq$ profitmaksimering, kvalitet kan ofres. [SLIDE]
\item \textbf{Expense center.} Fordel: praktisk når output ikke kan måles. Ulempe: subjektiv evaluering, empire building, soft budget constraints. [SLIDE]
\item \textbf{Revenue center.} Fordel: bevarer sentral priskontroll. Ulempe: omsetningsmaks ($MR=0$) gir for store rabatter / feil produktmiks. [SLIDE]
\item \textbf{Profit center.} Fordel: utnytter lokal spesifik viden. Ulempe: eksternaliteter mellom divisjoner, transfer pricing-konflikter, free-riding på fellesgoder. [SLIDE/BOK]
\item \textbf{Investment center.} Fordel: kobler drift og kapitalbruk. Ulempe: kortsiktige regnskapsmål (ROA) kan blokkere gode langsiktige investeringer; løsningen er EVA. [SLIDE/BOK]
\end{itemize}

\subsection*{Forutsetninger for valg av centertype}
\begin{itemize}\setlength{\itemsep}{2pt}
\item \textbf{Cost center:} output er standardisert og målbart; pris/marked styres sentralt; lokal kunnskap er begrenset til produksjonsmetode.
\item \textbf{Expense center:} output er ikke objektivt målbart; enheten leverer støttefunksjoner; budsjett + subjektiv vurdering er beste tilgjengelige kontrakt.
\item \textbf{Revenue center:} salgsinnsats er kontrollerbar lokalt, men toppledelsen vil eie prissetting (merkebeskyttelse, prisdiskriminering, refusjonsregimer).
\item \textbf{Profit center:} lederen sitter på spesifik viden om både kostnader og marked; divisjonen kan operere selvstendig; eksternaliteter mot andre enheter er små eller styrt via transfer pricing.
\item \textbf{Investment center:} alle profit center-forutsetninger oppfylt + investeringsbeslutninger krever lokal kunnskap; enheten er stor nok til at kapitalavkastning kan måles meningsfullt.
\end{itemize}

\subsection*{Real-world eksempler}

\textbf{Eks.~1 -- Equinor (NO).}\quad [EKS]
Letevirksomheten på norsk sokkel drives som et \emph{investment center}: leder bestemmer letebrønner og feltutbygginger og måles på avkastning på investert kapital. Mongstad-raffineriet drives historisk som et \emph{cost center} -- prisene settes globalt, så lederen måles på kostnadseffektivitet pr.\ raffinert fat. Eksemplet illustrerer hvordan \emph{samme} konsern bruker forskjellige centertyper avhengig av hvor lokal kunnskap og kontrollvariabler ligger.

\textbf{Eks.~2 -- Vy (tidl.\ NSB), sammenligning med Eks.~1.}\quad [EKS]
Persontrafikk på regulerte ruter ble lenge drevet som et klassisk \emph{cost center}: billetter og ruteplan var politisk fastsatt, så lederne kunne bare påvirke driftskostnadene. Da konsernet ble splittet i forretningsenheter, ble flere selskaper omgjort til \emph{profit centers} med ansvar for kommersielle ruter, bussvirksomhet og godstransport. Endringen \emph{endret atferd}: kommersielt fokus økte, men også konflikter om felles infrastruktur og merkenavn dukket opp -- klassisk profit center-problem. Sammenligning med Equinor viser samme prinsipp: centertypen følger hvor mye beslutningsrett som er reelt delegert.

\subsection*{Kobling til Mærsk Power-to-X}
I Mærsk er Ocean-divisjonen og Logistics \& Services drevet som separate \emph{profit/investment centers} med egne P\&L. Power-to-X-satsingen er typisk et stort investeringsprosjekt der ansvarsenheten må evalueres som et investment center (kapitalkrav er enormt og avkastningen kommer langt frem i tid). Caset belyser EVA-problematikken: regnskapsmessig ROA i de første årene kan se elendig ut, mens nåverdien er positiv. [CASE]

\subsection*{Fallgruber}
\begin{itemize}\setlength{\itemsep}{2pt}
\item Glemme controllability -- bruke profit som måltall der lederen ikke kontrollerer prisene.
\item Forveksle expense og cost center -- expense brukes når output \emph{ikke} kan måles.
\item Ikke nevne ROA/EVA-svakhetene ved investment centers (kortsiktig forvridning).
\item Hoppe over forutsetninger -- sensor vil ha både fordeler/ulemper \emph{og} valgkriterier.
\item Behandle centertyper isolert fra de tre søylene -- alltid koble til S1/S2/S3.
\end{itemize}

\subsection*{Håndskrivbar disposisjon (15-min forberedelse)}
\begin{itemize}\setlength{\itemsep}{2pt}
\item \textbf{Svar (kort):} Centertype = beslutningsrett (S1) + måltall (S2) + insentiv (S3) -- samtidig design. Valget følger controllability-prinsippet: måltallet skal dekke nøyaktig det lederen kan kontrollere.
\item \thref{responsibility-centers}{De 5 centertypene}: cost / expense / revenue / profit / investment -- stige av delegert kontroll.
\item Tabell: hvilken type kontrollerer kostnad / inntekt / profit / kapital.
\item Fordeler/ulemper pr.\ type -- ett poeng hver:
  \begin{itemize}
    \item Cost: enkelt, men kvalitet ofres.
    \item Expense: output uobserverbart, empire building.
    \item Revenue: $MR=0 \neq MR=MC$.
    \item Profit: spesifik viden, men eksternaliteter / transfer pricing.
    \item Investment: kapitalansvar, men ROA korttidsforvridd $\Rightarrow$ EVA.
  \end{itemize}
\item Forutsetning pr.\ type: hvor sitter lokal kunnskap, hva er målbart, hvilke priser settes sentralt.
\item \thref{controllability-principle}{Controllability}: målet skal dekke nøyaktig det lederen kan styre.
\item Eks.\ 1: Equinor (Leting = investment, Mongstad = cost). Eks.\ 2: Vy før/etter omorganisering.
\item Mærsk Power-to-X: investment center, EVA-poenget.
\item Søyle: \STwo\ primært; også \SOne\ og \SThree.
\end{itemize}
